% VTC2027-Spring Regular Paper.
\documentclass[conference]{IEEEtran}

\usepackage{amsmath,amssymb}
\usepackage{booktabs}
\usepackage{graphicx}
\usepackage{xcolor}
\usepackage{url}
\usepackage{stfloats}

\title{GPS L1 C/A Multipath Extraction and Conditional Modeling by Environment and Satellite Elevation Using SAGE}

\author{\IEEEauthorblockN{Liang Feng}
\IEEEauthorblockA{Tongji University, Shanghai, China\\2432132@tongji.edu.cn}}

\begin{document}
\maketitle

\begin{abstract}
GNSS multipath in dynamic road environments consists of propagation components that change as the vehicle moves. This paper extracts persistent path observations from real GPS L1 C/A measurements using navigation-data alignment, per-satellite SAGE estimation, and temporal consistency. A three-dimensional Gaussian mixture model jointly describes excess delay, the magnitude of relative Doppler, and relative power for Urban and Mountain/Valley measurements in three satellite elevation ranges. The fitted distributions reproduce the main delay and Doppler peaks and the two main relative-power peaks in the measurements. The model provides a joint description of these path parameters for each combination of environment and elevation range.
\end{abstract}

\begin{IEEEkeywords}
GNSS multipath, GPS L1 C/A, SAGE, Gaussian mixture model, environment and elevation modeling.
\end{IEEEkeywords}

\section{Introduction}
Global Navigation Satellite Systems (GNSS) support vehicle navigation and intelligent transportation. Reflections from buildings, terrain, vehicles, and other roadside structures change as the antenna moves. They produce propagation components with changing delay, Doppler, and power \cite{eissfeller1996gpsdynamic,mora1998multipath,beitler2015cmcd}. Carrier-to-noise-density ratio ($C/N_0$), pseudorange residuals, and positioning error describe the signal response at the receiver. Path-level extraction links this response to individual propagation components \cite{bilich2007snr,xie2011vehicular,beitler2015cmcd}.

Path-level analysis of real in-phase/quadrature (I/Q) measurements separates local signal components from components that persist over time. High-resolution parameter extraction provides the link between measured GNSS signals and propagation observations \cite{fleury1999sage}. We combine navigation-data-aided window formation, per-satellite SAGE estimation, temporal consistency, and conditional statistical modeling.

The paper has three main contributions. First, we present a multi-satellite GPS L1 C/A measurement and processing chain. Each tracked satellite is aligned and estimated separately. Second, we use temporal consistency to retain components that recur in adjacent windows. Third, we develop a three-dimensional Gaussian mixture model that preserves the joint behavior of excess delay, the magnitude of relative Doppler, and relative power across two environments and three elevation ranges.

\section{Measurement and Experimental Setup}
\subsection{Measurement Platform}
The measurement platform consists of a TEST-TREE RF-Catcher V2 for RF capture and playback and a GNSS dome antenna mounted on the vehicle roof. The antenna uses right-hand circular polarization (RHCP) and has a 40-dB active gain.

\subsection{Signal Acquisition Configuration}
The signal is GPS L1 C/A with a center frequency of 1575.42 MHz \cite{gpsis200n2022}. Samples are stored as interleaved in-phase and quadrature (I/Q) values in little-endian signed 16-bit format. The sampling rate is 10.23 MHz.

\subsection{Experimental Scenarios}
Two road-environment classes are considered: Urban and Mountain/Valley. Urban denotes roads surrounded by dense buildings. Mountain/Valley denotes roads in mountainous or valley terrain. Satellite elevation is divided into three ranges: LOW $=[0,30^\circ)$, MID $=[30^\circ,60^\circ)$, and HIGH $=[60^\circ,90^\circ]$.

\subsection{Processing Overview}
Fig.~\ref{fig:workflow} shows the measurement and processing chain. GNSS-SDR first tracks the signals and decodes the navigation data. We then process each tracked satellite separately. Navigation data provide signal timing and define valid 40-ms windows. We screen the windows, estimate delay and Doppler with SAGE, and retain components that remain consistent over time. Finally, we group the retained observations by environment and satellite elevation. GNSS-SDR follows the software-defined GNSS receiver architecture described in \cite{fernandez2011gnsssdr}.

\begin{table}[t]
\centering
\caption{Measurement and processing configuration.}
\label{tab:measurement-configuration}
\scriptsize
\setlength{\tabcolsep}{3pt}
\begin{tabular}{p{0.28\columnwidth}p{0.62\columnwidth}}
\toprule
Item & Configuration\\
\midrule
Acquisition device & TEST-TREE RF-Catcher V2\\
Antenna & GNSS dome antenna; RHCP; vehicle-roof mounting\\
Signal & GPS L1 C/A; center frequency 1575.42 MHz\\
Sampling & 10.23 MHz (10230000 Hz)\\
IQ format & Interleaved I/Q; little-endian signed int16\\
GNSS-SDR processing & Tracking and navigation decoding\\
Satellite processing & Each tracked satellite processed separately before aggregation\\
Path processing & Navigation alignment, candidate screening, SAGE estimation, temporal consistency, and modeling by environment and elevation\\
\bottomrule
\end{tabular}
\end{table}

\begin{figure*}[t]
\centering
\includegraphics[width=\textwidth]{../../figures/figure1_multisatellite.pdf}
\caption{Multi-satellite measurement and processing workflow. Each tracked satellite passes through navigation alignment, local SAGE estimation, and a time-consistency check. The resulting path observations are then weighted and grouped by environment and satellite elevation.}
\label{fig:workflow}
\end{figure*}

\section{Navigation-Aided SAGE and Conditional Path Model}
\subsection{Per-Satellite Signal and Multipath Model}
Let $M$ be the number of tracked satellites. For satellite $m$, the complex baseband signal is modeled as
\begin{equation}
r_m(t)=\sum_{\ell=0}^{L_m-1}\alpha_{m,\ell}s_m(t-\tau_{m,\ell})\exp\!\left(j2\pi\Delta f_{m,\ell}t\right)+n_m(t),
\label{eq:multisat-model}
\end{equation}
Here, $s_m(t)$ is the navigation-aligned code waveform and $n_m(t)$ is noise. The satellite index is $m=1,\ldots,M$. The value $L_m$ gives the number of components in the local model. The index $\ell=0$ denotes the direct component. The direct-path delay $\tau_{m,0}$ is the reference delay for satellite $m$. Relative Doppler is measured with respect to this component, so $\Delta f_{m,0}=0$. Each component has complex gain $\alpha_{m,\ell}$, delay $\tau_{m,\ell}$, and relative Doppler offset $\Delta f_{m,\ell}$. The complex gain contains amplitude and phase.

\subsection{Navigation-Data-Aided Observation Formation}
Tracking results and decoded navigation data provide the timing needed to select samples. The decoded bits identify the satellite's pseudorandom-noise (PRN) code and provide a known sign sequence aligned in time. Removing this sequence enables coherent processing over complete 40-ms windows. The PRN-to-channel mapping selects the desired tracked signal. At 10.23 MHz, each C/A chip contains 10 samples.

\subsection{Per-Satellite SAGE Delay--Doppler Estimation}
We use the space-alternating generalized expectation-maximization (SAGE) algorithm for high-resolution delay--Doppler estimation \cite{fessler1994sage,fleury1999sage}. At iteration $i$, we isolate component $\ell$ for satellite $m$ by subtracting the synthesized contributions of the other components,
\begin{equation}
r_{m,\ell}^{(i)}(t)=r_m(t)-\sum_{k\ne\ell}\hat{\alpha}_{m,k}^{(i)}q_m(t;\hat{\tau}_{m,k}^{(i)},\hat{\Delta f}_{m,k}^{(i)}),
\label{eq:hidden-signal}
\end{equation}
where $q_m(\cdot)$ is the navigation-aligned code replica with fractional delay and Doppler. The current delay and Doppler are refined by maximizing
\begin{equation}
(\hat{\tau}_{m,\ell},\hat{\Delta f}_{m,\ell})=\mathop{\arg\max}_{\tau,\Delta f}
\frac{|q_m(\tau,\Delta f)^{H}r_{m,\ell}^{(i)}|^2}{q_m(\tau,\Delta f)^{H}q_m(\tau,\Delta f)}.
\label{eq:sage-objective}
\end{equation}
First, we compute coarse delay--Doppler correlation with the fast Fourier transform (FFT) and its inverse (IFFT). We then refine the result with fractional-delay replicas. Complex gains are updated by least squares. Iterations stop after at most 10 updates or when the relative change in residual sum of squares (RSS) is below $10^{-6}$. The local delay grid uses 0.1-sample spacing (0.01 chip) over a 0.8-sample neighborhood. The local Doppler grid spans $\pm30$ Hz at 5-Hz spacing around the tracking estimate.

For each screened window, we evaluate models with $L_m=1,2,3,4$. A higher-order model is used when the fit is numerically valid, lowers the Bayesian information criterion (BIC) by at least 10, and reduces the residual sum of squares by at least 0.002\%. We then use temporal consistency to retain components that recur over time.

\subsection{Persistent Path Observations}
After local estimation, we compare path parameters in neighboring windows. We retain a component when matching components occur in at least three consecutive windows within a two-window neighborhood. Matching uses tolerances of 1.5 samples, 40 Hz, and 10 dB for delay, relative Doppler, and relative power, respectively. The center window of each sequence defines one persistent path observation. Each observation includes its satellite, recording, measurement scene, observation time, and associated geometry.

Using the direct component as the reference, we define the three modeled parameters as
\begin{equation}
\begin{aligned}
\Delta\tau_{m,\ell}&=\tau_{m,\ell}-\tau_{m,0},\\
f^{\mathrm{rel}}_{m,\ell}&=\Delta f_{m,\ell},\\
P_{m,\ell}&=10\log_{10}\!\frac{|\alpha_{m,\ell}|^2}{|\alpha_{m,0}|^2}.
\end{aligned}
\label{eq:path-parameters}
\end{equation}
Here, $\Delta\tau$ is excess delay in samples, $f^{\mathrm{rel}}$ is signed Doppler relative to the direct component in Hz, and $P$ is relative path power in dB.

Window-level geometry alignment provides a satellite elevation for each path observation. Observations with valid alignment enter the elevation groups, and the full observation set is used for the environment-level summaries.

\subsection{Conditional Three-Dimensional Mixture Model}
Let $e_i$ and $b_i$ denote the environment and elevation range of observation $i$. Each observation has weight $w_i=1/N_{\mathrm{track}(i)}$. Thus, every satellite track contributes the same total weight. For an environment and elevation group $(e,b)$, the weighted empirical distribution of parameter $p$ is
\begin{equation}
\widehat F_{e,b,p}(x)=
\frac{\sum_{i\in\mathcal I_{e,b}}w_i\,\mathbf{1}(z_{i,p}\le x)}
{\sum_{i\in\mathcal I_{e,b}}w_i}.
\label{eq:weighted-cdf}
\end{equation}
where $z_{i,p}$ is the value of parameter $p$ for observation $i$.

We transform the three parameters into
\begin{equation}
\mathbf{x}_i=\left[\log \Delta\tau_i,\ \log\!\left(1+|f_i^{\mathrm{rel}}|/1\,\mathrm{Hz}\right),\ P_i\right]^{\mathsf T}
\label{eq:gmm-features}
\end{equation}
and standardize each feature using weighted means and standard deviations from the training fold. The Doppler feature uses the magnitude of relative Doppler to represent the separation between paths.

For environment $e$ and elevation range $b$, the conditional density is
\begin{equation}
p(\mathbf{x}\mid e,b)=\sum_{k=1}^{3}\pi_{e,b,k}\,
\mathcal N\!\left(\mathbf{x};\boldsymbol\mu_{e,k},\boldsymbol\Sigma_k\right).
\label{eq:conditional-gmm}
\end{equation}
The component means depend on the environment. The mixture weights depend on both environment and elevation, while the covariance matrices are shared across groups. The model uses three Gaussian components. We selected this number by holding out one measurement scene at a time during validation.

\section{Experimental Results}
\begin{figure*}[!t]
\centering
\includegraphics[width=\textwidth]{../../figures/conditional_empirical_vs_model_pdf_environment_elevation.pdf}
\caption{Track-weighted measured distributions (step dashed lines) and fitted marginal probability densities (solid lines) for the six combinations of environment and elevation range.}
\label{fig:empirical-model-pdf}
\end{figure*}

\subsection{Modeling Data by Environment and Satellite Elevation}
The model uses 518 persistent path observations from 236 satellite tracks, 36 measurement recordings, and 9 measurement scenes. Window-level geometry alignment provides valid elevation values for 487 observations. These 487 observations provide the six combinations of environment and elevation range. All 518 observations contribute to the environment-level summaries. Table~\ref{tab:environment-elevation-cells} lists the six modeled groups. Here, $N_{\mathrm{eff}}$ is the effective observation count after track balancing.

\begin{table*}[t]
\centering
\caption{Path observations by environment and elevation range.}
\label{tab:environment-elevation-cells}
\scriptsize
\setlength{\tabcolsep}{5pt}
\begin{tabular}{llrrrrr}
\toprule
Environment & Elevation range & Observations & Tracks & Recordings & Scenes & $N_{\mathrm{eff}}$\\
\midrule
Urban & LOW $[0,30^\circ)$ & 18 & 8 & 3 & 3 & 14.8\\
Urban & MID $[30^\circ,60^\circ)$ & 169 & 74 & 10 & 5 & 125.6\\
Urban & HIGH $[60^\circ,90^\circ]$ & 129 & 64 & 8 & 5 & 101.1\\
Mountain/Valley & LOW $[0,30^\circ)$ & 22 & 12 & 3 & 3 & 17.3\\
Mountain/Valley & MID $[30^\circ,60^\circ)$ & 117 & 53 & 8 & 3 & 86.9\\
Mountain/Valley & HIGH $[60^\circ,90^\circ]$ & 32 & 15 & 3 & 2 & 23.1\\
\bottomrule
\end{tabular}
\end{table*}

\subsection{Measured and Fitted Parameter Distributions}
Figure~\ref{fig:empirical-model-pdf} compares the track-weighted measured histograms with the marginal densities from the joint model. All six groups have a main peak near 1.3--1.6 samples and a tail toward larger delays. The fitted curves reproduce the main peak and the overall tail while smoothing local fluctuations.

Figure~\ref{fig:joint-delay-power} shows a two-dimensional delay--power view of the joint model. Blue points are measured observations. Orange solid and dashed contours enclose the central 50\% and 90\% density regions, respectively. The two environments show two main combinations: shorter delays with stronger power and longer delays with weaker power.

\begin{figure}[!t]
\centering
\includegraphics[width=0.96\columnwidth]{../../figures/conditional_joint_delay_power.pdf}
\caption{Measured and fitted joint delay--power density.}
\label{fig:joint-delay-power}
\end{figure}

\clearpage

The magnitude of relative Doppler has a clear peak near 45--50 Hz in all six combinations of environment and elevation range. The fitted curves reproduce this peak and show the secondary structure and the tail at higher Doppler values in smooth form. Similar peak locations across the elevation ranges show a common dominant Doppler characteristic.

Relative power has two recurring concentration regions near $-14$ and $-4$ dB. The three-component model represents this bimodal structure. In Urban measurements, the relative weight of the two regions changes from LOW to HIGH elevation. Mountain/Valley shows both regions in all three elevation ranges, with a different balance between them. The model captures the common two-region structure and the group-dependent balance.

Together, these plots show the main marginal and joint structure of the three path parameters.

\section{Conclusion}
We present a navigation-aided processing chain for extracting persistent GPS L1 C/A path observations from dynamic road measurements. Each tracked satellite is aligned and processed separately with fractional delay--Doppler SAGE. Temporal consistency is then applied before track-weighted aggregation.

A three-dimensional GMM jointly models excess delay, the magnitude of relative Doppler, and relative power for two environments and three elevation ranges. The measured and fitted marginal distributions, together with the joint delay--power plot, show the main delay and Doppler peaks, the two-region power structure, and the principal delay--power combinations. Together, these results give a common statistical description of the three path parameters for all six combinations of environment and elevation range.

\bibliographystyle{IEEEtran}
\bibliography{references}

\end{document}
